\documentclass{article}
\usepackage[utf8]{inputenc}
\usepackage[T1]{fontenc}
\usepackage{spconf,amsmath,graphicx,xcolor,hyperref}
\hypersetup{hidelinks,hypertexnames=false}
\usepackage{titlesec}
\usepackage{float}
\usepackage{booktabs}
\usepackage{balance}


\ninept

% \titleformat{\section}[block]
%   {\centering\normalsize\bfseries}
%   {\thesection.}
%   {1em}
%   {\MakeUppercase}

% \titlespacing*{\section}
%   {0pt}
%   {0.9ex}
%   {0.45ex}
\newcommand{\shap}{\ensuremath{\phi}}

\title{DO SIMPLEX FUSION WEIGHTS MEASURE REPRESENTATION UTILITY?\\
MOS ATTRIBUTION AND SYSTEM-DISJOINT SELECTION}

\name{\shortstack{Daniel Casanova$^{1,3}$ \quad Alef Iury Ferreira$^{1,2}$ \quad
Lucas Gris$^{1,2}$ \\
Pedro Lustosa Rege Botelho$^{1}$ \quad Fernanda Silva$^{1}$ \quad
Frederico Oliveira$^{1,2}$ \\
Arlindo Galvão Filho$^{1,2}$ \quad Anderson da Silva Soares$^{1,2}$}}

\address{\shortstack{
$^1$ AKCIT, Brazil \\
$^2$ Federal University of Goiás (UFG), Brazil \\
$^3$ Federal University of Technology -- Paraná (UTFPR), Brazil}}

\begin{document}
\ninept
\maketitle


\begin{abstract}
MOS predictors increasingly combine pretrained speech representations, but fusion weights describe allocation under a model constraint, not necessarily predictive utility. We analyze ten frozen ridge experts on four benchmarks with a non-negative simplex gate, exact Shapley attribution over all coalitions, and nested system-disjoint fixed-budget selection. Under simplex fusion, 16/40 corpus--representation pairs had near-zero full-gate weight and positive validation Shapley utility; 15/16 remained positive on test after the discovery list was frozen. An exploratory positive-NNLS audit found positive test Shapley values in 8/9 zero-weight cases across BVCC, SingMOS, and TMHINT-QI, but 0/4 for BRSpeechMOS, where NNLS overfit. In held-system selection, Shapley showed no multiplicity-adjusted advantage over exhaustive or singleton selection. Excluding BRSpeechMOS reversed the Shapley--exhaustive point estimate from $-2.15\%$ to $+0.63\%$. Thus, the primary allocation--utility finding is specific to simplex fusion, and compact systems require direct system-disjoint validation.
\end{abstract}



\begin{keywords}
mean opinion score, speech quality, representation analysis, Shapley value,
self-supervised learning
\end{keywords}

\section{Introduction}
\label{sec:intro}

Subjective listening tests remain the reference standard for assessing
synthetic and processed speech, but they are costly and difficult to scale.
Non-intrusive Mean Opinion Score (MOS) prediction offers an automatic
alternative by learning to reproduce human quality ratings. Early systems such
as MOSNet used task-specific convolutional and recurrent architectures and
reported substantially stronger correlations at the system level than at the
utterance level~\cite{lo2019mosnet}. More recent systems use pretrained
self-supervised speech representations, which improved performance on
benchmarks such as VCC2018~\cite{tseng2021sslmos}.

Current MOS predictors increasingly combine several representations. Prior work
has fused multiple encoders from the wav2vec, HuBERT, and WavLM families and
reported improvements over individual predictors, including under
out-of-domain conditions~\cite{yang2022fusion}. Other systems supplement
acoustic representations with pitch, duration, TTS-encoder states, linguistic
tags, or text embeddings~\cite{vioni2023contentaware}. These results establish
that heterogeneous representations can be complementary. They do not,
however, establish that the magnitude of a learned fusion coefficient measures
the general importance of its representation. The coefficient depends on
scaling, regularization, correlation, model constraints, and the alternatives
available in the mixture. In a different setting, Jain and Wallace likewise
caution that attention weights alone need not explain model predictions
\cite{jain2019attention}.

Generalization is a second challenge. Strong results within one benchmark do
not necessarily transfer to unseen synthesis systems or assessment conditions
\cite{cooper2022generalization,huang2022voicemos}. This distinction also affects
experimental design. A random utterance-level split may place samples from the
same synthesis system in both training and evaluation, measuring interpolation
within known systems rather than transfer to a new one. When the intended claim
concerns unseen systems, the evaluation must preserve system identity across
the split~\cite{roberts2017crossvalidation}.

The interpretation of a fused model and the selection of a compact model are
therefore related but different problems. A gate weight measures allocation in
the complete mixture. Singleton loss measures an expert in isolation, whereas
leave-one-out loss measures its necessity after every other expert is present.
Shapley utility instead averages the expert's marginal contribution across
coalitions~\cite{shapley1953value,covert2020sage}. These measures can disagree
when representations encode correlated or redundant information
\cite{kumar2020shapley}. A representation may receive little weight in the
complete mixture while remaining useful when a stronger substitute is absent.
This issue is especially relevant here because each player represents a
complete checkpoint, pooling rule, preprocessing path, and linear readout
rather than an isolated scalar feature. Moreover, Shapley utility is
exponential in the number of players and does not directly optimize any
fixed-size subset.

In this paper, we study ten frozen representation pipelines on BRSpeechMOS,
BVCC, SingMOS, and TMHINT-QI, with a non-negative simplex gate over linear
ridge experts as the primary fusion model. We ask whether a representation can
receive near-zero weight in that full gate while retaining predictive utility
in other coalitions, and whether coalition-based attribution can identify
compact subsets that generalize to unseen systems. We answer the first question
with the complete ten-player simplex lattice, supplemented by an exploratory
positive-NNLS sensitivity, and the second through nested system-disjoint
selection with equal-system MSE.


Our aim is not to introduce Shapley attribution as a new measure of representation importance. Rather, we use it to examine whether three notions that are often treated as interchangeable actually agree: allocation in the full mixture, average contribution across alternative coalitions, and the value of a representation for selecting a compact model that generalizes to unseen systems. The primary coalition analysis uses simplex fusion; the NNLS analysis is a secondary, test-set sensitivity. The nested system-disjoint benchmark evaluates selection separately. This separation allows us to test whether contextual contribution, although informative about redundancy and substitutability, is also useful for prospective representation selection.

The contributions of this work are as follows:
\begin{itemize}
  \setlength{\itemsep}{0pt}
  \setlength{\parsep}{0pt}
  \setlength{\topsep}{2pt}
  \item We introduce a validation-frozen audit for the simplex fusion of
  frozen linear experts, separating full-mixture allocation from
  coalition-averaged predictive utility.
  \item Under primary simplex fusion, we identify 16 of 40
  corpus--representation pairs with near-zero full-gate weight but positive
  point-estimate Shapley utility, and test sensitivity to refits and player
  definitions; a separate NNLS analysis is exploratory.
  \item We compare gate, singleton, Shapley, and exhaustive selection in a
  nested held-system benchmark. Exact Shapley attribution diagnoses contextual
  redundancy but provides no multiplicity-adjusted selection advantage over the
  ten-score singleton baseline.
\end{itemize}

\section{Methods}
\label{sec:methods}

Our experimental framework has two stages. We first perform a complete-coalition
audit to compare full-model allocation with contextual utility. We then test
fixed-budget selection using nested system-disjoint cross-validation.

\noindent\textbf{Data and representations.}
We evaluate BRSpeechMOS~\cite{oliveira2023representations}, BVCC~\cite{huang2022voicemos}, SingMOS~\cite{tang2024singmos}, and TMHINT-QI~\cite{chen2022inqss} independently using their official train/validation/test partitions. In corpus order, these partitions contain 1{,}942/242/243, 4{,}974/1{,}066/1{,}066, 2{,}000/544/645, and 11{,}644/1{,}293/1{,}978 utterances; the training partitions contain 6, 175, 31, and 98 systems. Audio is loaded with torchaudio, downmixed by channel averaging, and resampled with \nolinkurl{torchaudio.functional.resample} to 16~kHz. System identity defines every grouped fold.

The ten frozen inputs are \nolinkurl{openai/whisper-large-v3} (last hidden state, 1,280-D),
\nolinkurl{lengyue233/content-vec-best} (layer 12, 768-D),
\nolinkurl{microsoft/wavlm-base-plus} (layer 12, 768-D),
\nolinkurl{Bencr/beats-checkpoints/BEATs_iter3_plus_AS2M.pt} (final hidden state, 768-D),
\nolinkurl{mispeech/ced-small} (final encoder state, 384-D),
\nolinkurl{speechbrain/spkrec-ecapa-voxceleb} (192-D, $\ell_2$-normalized),
a 64-band ERB frontend, continuous and quantized RMVPE, and eGeMAPSv02
\cite{radford2023whisper,qian2022contentvec,chen2022wavlm,chen2023beats,desplanques2020ecapa,wei2023rmvpe,dinkel2024ced,eyben2016egemaps}.
Encoder frame sequences are mean-pooled; Whisper extraction uses the last
hidden layer with padding cropping enabled. ERB uses a power STFT (512-point
FFT, 400-sample Hann window, 160-sample hop), triangular filters from
50--8{,}000~Hz, $\log(1+x)$ compression, and temporal averaging. RMVPE uses
the \nolinkurl{weights/rmvpe.pt} checkpoint from RVC-Boss/RMVPE, with a
160-sample hop at 16~kHz and unvoiced threshold 0.03 (unvoiced frames have
$F_0=0$); continuous features are
voiced-$F_0$ mean, population standard deviation, minimum, maximum, and voiced
fraction. Quantized RMVPE is a normalized 256-bin utterance histogram (bin 0
unvoiced; 255 mel-spaced voiced bins spanning 50--1{,}100~Hz). eGeMAPSv02
comprises 88 openSMILE functionals. MOS labels do not adapt the encoders.
Checkpoint family identifiers and layers were recorded, but hub revisions,
weight-file hashes, and the torchaudio/openSMILE package versions were not;
bit-identical re-extraction is therefore not claimed.

Each representation $j$ feeds the same ridge probe ($\alpha=10$, unpenalized intercept; scikit-learn LSQR solver, tolerance $10^{-10}$, maximum 10{,}000 iterations). Feature means and standard deviations are estimated within each training fold and reused unchanged on its scoring fold; dimensions with standard deviation below $10^{-6}$ use scale one. eGeMAPS is additionally clipped at training-fold 0.1/99.9 percentiles. This fold-wise preprocessing prevents information from the scoring systems entering feature construction.


\noindent\textbf{Coalition audit.}
We use sequential least-squares quadratic programming (SLSQP) to combine the
experts with a non-negative simplex gate:
\begin{equation}
\hat y_S(x)
=
\sum_{j\in S} w_j^{(S)} f_j(x),
\qquad
w_j^{(S)}\geq 0,
\qquad
\sum_{j\in S} w_j^{(S)}=1 .
\label{eq:simplex}
\end{equation}

Here, $S$ is the active coalition and $f_j(x)$ is expert $j$; experts retain
their unpenalized intercepts, and the fusion has no additional intercept.
Five system-grouped folds produce cross-fitted expert predictions, with all
preprocessing and ridge fitting restricted to each fold's training systems.
For every non-empty $S\subseteq N$, where $N$ contains the ten pipelines,
expert predictions remain fixed and only the simplex gate is fitted. Gate
weights therefore measure allocation within a fitted mixture.

We minimize training-OOF MSE with SLSQP, an analytic gradient, uniform
initialization, bounds $[0,1]$, the sum-to-one constraint, tolerance
$\texttt{ftol}=10^{-12}$, and at most 1{,}000 iterations. If SLSQP reports
failure, a deterministic projected-gradient fallback runs for 4{,}000 steps
from the uniform point with step size $0.2/\max(\|P\|_2^2,1)$, then projects
onto the simplex and renormalizes the weights. The same
solver is used for equal-system gates after scaling each system's rows by
$1/\sqrt{n_g}$. Saved gate vectors are finite and simplex-feasible, but
per-coalition optimizer status and fallback use were not retained; we therefore
do not report an all-coalition convergence rate.

We define $v(S)=-\mathcal{L}(S)$; the empty coalition predicts the MOS mean
estimated on the corresponding training partition and applied unchanged to the
scoring split. Higher utility thus means lower error. The audit uses
utterance-level MSE. Validation labels identify the audit cases, which are
frozen before test evaluation. Official-test labels only score frozen
predictors and compute post-hoc attribution quantities; they do not fit
experts, gates, preprocessing, coalition models, or selection rules.

The exact Shapley utility~\cite{shapley1953value} is
\begin{equation}
\phi_j
=
\sum_{S\subseteq N\setminus\{j\}}
\frac{|S|!\,(|N|-|S|-1)!}{|N|!}
\left[
v(S\cup\{j\})-v(S)
\right].
\label{eq:shapley}
\end{equation}

All $2^{10}-1=1{,}023$ non-empty coalitions are evaluated exactly. The gate is
fitted on cross-fitted training predictions and then frozen on the scoring
split. Leave-one-out utility measures the loss change after removing one
expert from the full coalition; Shapley utility measures contextual predictive
contribution, not causal effect. We define a gate as numerically zero when
$w_j<10^{-6}$.

As an exploratory fusion-rule sensitivity, we also recomputed the 1{,}023
coalition values for uniform averaging, positive NNLS with an intercept and no
sum-to-one constraint, and ridge stacking selected by grouped cross-validation
over 17 log-spaced $\alpha$ values from $10^{-4}$ to $10^4$. These analyses
reuse the frozen ridge experts and use test utterance-MSE; unlike the primary
validation-discovery audit, they are descriptive test-set comparisons without
bootstrap intervals.

Sensitivity analyses vary this threshold from $10^{-8}$ to $10^{-3}$, resample
training systems for 100 gate refits and 50 complete-game refits, and alter the
player definition by merging, grouping, or removing overlapping prosodic
pipelines. A quadratic assignment procedure (QAP) with 100{,}000 node-label
permutations relates prediction correlation to pairwise Shapley interaction;
accuracy-matched Gaussian and full-refit controls test whether prediction
geometry reproduces the pattern~\cite{grabisch1999interaction,krackhardt1988qap}.
These analyses are exploratory robustness diagnostics.

\noindent\textbf{System-disjoint selection.}
We evaluate subset selection separately on each corpus's training partition
using nested system-disjoint cross-validation. Each corpus has five outer
GroupKFold folds by training-system identifier; the deterministic partition is
not shuffled or repeated. Within each outer fold, all preprocessing, ridge fitting, fusion
calibration, attribution, and selection use only outer-development systems;
four grouped inner folds select subsets for $k\in\{1,2,3,5\}$.

We compare four rules. \emph{Gate}, \emph{singleton}, and \emph{Shapley} rank
representations by full-gate weight, individual predictive performance, and
exact inner-loop Shapley utility, respectively. \emph{Exhaustive} directly
selects the size-$k$ subset with the lowest inner loss. After selection,
preprocessing and candidate experts are refit on all outer-development systems;
the selected subset and calibrated gate are then frozen and evaluated once on
outer-held systems.

Both inner selection and outer evaluation use equal-system MSE,


\begin{equation}
\mathcal{L}_{\mathrm{sys}}
=
\frac{1}{|\mathcal{G}|}
\sum_{g\in\mathcal{G}}
\frac{1}{n_g}
\sum_{i:g_i=g}
(y_i-\hat y_i)^2 ,
\label{eq:systemmse}
\end{equation}
where $\mathcal{G}$ is the set of held systems, $n_g$ is the number of
utterances from system $g$, and $g_i$ identifies utterance $i$'s system. Each
system therefore contributes equally. Pairwise contrasts use 10{,}000 paired
bootstrap replicates that resample outer-held systems while preserving pairing
across methods and budgets. Holm's procedure corrects five pre-specified
contrasts: each ranking rule versus exhaustive search, gate versus singleton,
and Shapley versus singleton~\cite{holm1979procedure}.
The nested bootstrap uses seed 20260916; sensitivity reaggregations use
20260917 plus a scenario index. The primary attribution protocol and its
all-cell test bootstrap use seeds 20260910 and 20260909, respectively. The
validation-discovery confirmation of the frozen 16 cases uses 10{,}000 test-
system resamples, with seed $20260916+1000d+j$, where $d$ is the zero-based
corpus index in the order above and $j$ is the zero-based representation
index. Training-system refit draws use seed $20260919+1009d+104729r$ for
zero-based corpus index $d$ and replicate $r$; matched-null controls use seed
20260912. The QAP uses 100{,}000 label
permutations per corpus, seeded by
$20270910+\sum_c\mathrm{ord}(c)$ over the lowercase corpus identifier. The
separate SLSQP multistart audit uses seed 20260915 (offset by corpus index).
Because the grouped folds are deterministic, bootstrap replicates quantify
system-resampling uncertainty conditional on one fold assignment; they are not
repeated group splits.
The recorded environment was Python 3.10.12, NumPy 2.2.6, SciPy 1.15.3,
pandas 2.3.3, scikit-learn 1.7.2, and PyTorch 2.5.1+cu118; all reported
analyses ran CPU-only.

% Here, $\mathcal{G}$ is the set of systems, $n_g$ is the number of utterances
% from system $g$, $g_i$ is the system associated with utterance $i$, and
% $\hat y_i$ is the final fused prediction. Thus, each system contributes one
% equally weighted MSE regardless of its number of utterances.





\section{Results and Discussion}
\label{sec:results}

\textbf{Selection performance on unseen systems.}
We first compare the four selection rules on outer-held systems. The nested
benchmark uses equal-system MSE for both inner selection and outer evaluation.
Table~\ref{tab:selection_macro} reports the macro contrasts, and
Fig.~\ref{fig:selection} reports the corpus--budget cells. Negative values
indicate lower held-system error than the exhaustive inner-development
reference. Shapley has the lowest point estimate ($-2.15\%$, 95\% CI
$[-3.77,0.71]\%$), followed by gate ($-1.51\%$, $[-2.82,0.74]\%$) and
singleton ($-0.72\%$, $[-2.91,1.44]\%$). No selector showed
multiplicity-adjusted evidence of a difference. Percentages are normalized by
the full-model held-system MSE, and macro values average the 16
 corpus--budget cells.

Because BRSpeechMOS contributes only six training systems, we reaggregated the
saved held-system results after excluding it (three corpora and 12
corpus--budget cells; no models were refit).
Relative to exhaustive selection, the deltas were $+0.76\%$
[$+0.04,+1.54$] for gate ($p_{\mathrm{Holm}}=.156$), $+0.78\%$
[$+0.14,+1.46$] for singleton ($p_{\mathrm{Holm}}=.094$), and $+0.63\%$
[$-0.02,+1.30$] for Shapley ($p_{\mathrm{Holm}}=.172$); brackets are
unadjusted 95\% paired system-bootstrap CIs. Thus, excluding BRSpeechMOS
reverses the all-corpus point estimates, including Shapley's change from
$-2.15\%$ to $+0.63\%$. No contrast survives Holm correction in the
all-corpus or BRSpeechMOS-excluded comparison. The inferential conclusion is
unchanged, but the point-estimate
ranking is sensitive to the small BRSpeechMOS corpus; we therefore do not
interpret the all-corpus ranking as corpus-robust.

We also performed leave-one-corpus-out macro reaggregation of the same
outer-held errors (not model refitting or cross-corpus transfer). Shapley
versus exhaustive deltas were $-3.49\%$ [$-5.57,+0.29$] without BVCC,
$-2.93\%$ [$-5.05,+0.85$] without SingMOS, and $-2.79\%$
[$-4.94,+0.99$] without TMHINT-QI; corresponding within-scenario Holm
$p$-values were $.494$, $.706$, and $.673$; none of these Shapley--exhaustive
contrasts survived correction. The displayed leave-one-corpus-out intervals
are unadjusted 95\% paired system-bootstrap CIs; Holm correction is within
each omission scenario and does not adjust across the four scenarios.
One exploratory Shapley--singleton contrast was significant only when omitting
SingMOS ($-1.82\%$, 95\% CI $[-4.18,-0.14]\%$,
$p_{\mathrm{Holm}}=.049$ within that scenario); it was not corrected across
the four leave-one-corpus-out scenarios and is not treated as confirmatory.

\begin{table}[t]
\centering
% \caption{Macro held-system comparison with exhaustive inner-development
% selection as the reference. Equal-system MSE is used for both inner selection
% and outer evaluation. Values are averaged over the four corpora and primary
% budgets $k\in\{1,2,3,5\}$; negative values indicate lower held-system MSE than
% the reference. Intervals are paired 95\% system-bootstrap intervals, and
% $p$-values are Holm-adjusted. The score column counts coalition evaluations,
% not runtime or encoder cost. Percentages are normalized by the full-model
% held-system MSE; the macro average covers 16 corpus--budget cells.}
\caption{Macro held-system comparison with exhaustive inner-development
selection as reference. Values average four corpora and
$k\in\{1,2,3,5\}$; negative values indicate lower held-system MSE.
Intervals are paired 95\% system-bootstrap CIs and $p$-values are
Holm-adjusted. Scores count coalition evaluations, and percentages are
relative to full-model held-system MSE.}

\label{tab:selection_macro}
\setlength{\tabcolsep}{3.5pt}
\begin{tabular}{lrrc}
\toprule
Selector & \shortstack{Ranking/coalition\\scores} & $\Delta$MSE [95\% CI] & $p_{\mathrm{Holm}}$ \\
\midrule
Gate      & 1       & $-1.51$ [$-2.82,+0.74$]   & $.694$ \\
Singleton & 10      & $-0.72$ [$-2.91,+1.44$]   & $1.000$ \\
Shapley   & 1{,}023 & $-2.15$ [$-3.77,+0.71$]   & $.694$ \\
\bottomrule
\end{tabular}
\end{table}

\begin{figure*}[t]
\centering
\includegraphics[width=0.97\textwidth]{fig_selection_story.pdf}
\caption{Held-system equal-system MSE change by corpus and subset budget,
relative to exhaustive inner-development selection. Values are normalized by
full-model MSE (negative is better); asterisks mark cellwise 95\% system-
bootstrap intervals excluding zero, without multiplicity adjustment. Macro
intervals and Holm-adjusted contrasts are in Table~\ref{tab:selection_macro}.}
\label{fig:selection}
\end{figure*}

% Across the five prespecified contrasts, no multiplicity-adjusted comparison
% showed evidence of a difference. The Shapley--singleton contrast favored
% Shapley by $1.43\%$ at the point estimate, with a 95\% CI of
% $[-3.23,-0.07]\%$, but its raw $p=.0296$ became
% $p_{\mathrm{Holm}}=.148$. This result does not establish equivalence between
% the selectors.

The Shapley-minus-singleton contrast was $-1.43\%$ (95\% CI
$[-3.23,-0.07]\%$; raw $p=.0296$), but it was not significant after Holm
correction ($p_{\mathrm{Holm}}=.148$). This absence of adjusted significance
does not establish equivalence: no equivalence margin or equivalence test was
pre-specified.

The selectors also retained different subsets. At $k=1$, singleton and
exhaustive search coincide by construction, and all selectors retain CED on
BVCC in every outer fold. At $k=5$ on TMHINT-QI, all rules retain WavLM and
BEATs; CED is retained in three to five folds depending on the selector. On
BRSpeechMOS at $k=5$, eGeMAPS is selected in all five outer folds by gate,
singleton, Shapley, and exhaustive selection. SingMOS reverses direction
between $k=3$ and $k=5$. Similar held-system errors can therefore arise from
different subsets.

Singleton ranking measures standalone predictive quality, whereas Shapley
averages marginal gains across coalitions of different sizes. Deployment
evaluates a particular fixed-budget subset. A representation can contribute
when a stronger alternative is absent while remaining unnecessary in the
selected subset. Exact Shapley estimation required 1{,}023 coalition scores,
compared with ten representation-specific scores for singleton ranking. These
counts measure scoring requirements rather than encoder latency or energy;
broader attribution did not yield a detectable held-system selection gain in
this experiment.

\textbf{Simplex-fusion allocation does not determine contextual utility.}

In the primary non-negative simplex audit, we compare full-gate allocation with contextual utility. Using the
operational threshold $w_j<10^{-6}$, the validation-discovery audit
identified 16 of 40 corpus--representation pairs with numerically zero
full-gate weight and positive point-estimate Shapley utility
(Fig.~\ref{fig:attribution}). The count remained 16 as the threshold varied
from $10^{-8}$ to $10^{-3}$, although exact equality to zero remains
solver-tolerance dependent. On the frozen test evaluation, 15/16 cases
remained positive at the point estimate, 11/16 had an unadjusted
system-bootstrap interval with a lower bound above zero, and 8/16 were
positive in all 50 training-system refits. Five effects were below 2\% of
full-model MSE, whereas BRSpeechMOS Speaker/CED and BVCC/TMHINT-QI eGeMAPS
exceeded 17\%. BRSpeechMOS estimates remain exploratory because its training
partition contains only six systems.

\begin{figure*}[t]
\centering
\includegraphics[width=0.97\textwidth]{fig_attribution_utility.pdf}
\caption{Full-gate weight and test Shapley utility for all 40
corpus--representation pairs (a), and test-system 95\% bootstrap intervals
for the 16 cases frozen by validation (b). Crosses mark those cases; near-zero
weights are display-jittered. Colors identify the four corpora, and filled
markers have a positive interval lower bound.}
\label{fig:attribution}
\end{figure*}

All 16 discordant cases had absolute leave-one-out utility below
$1.1\times10^{-7}$ MSE. Removing the corresponding representation after all
other experts were present therefore had negligible effect. Positive Shapley
utility indicates a positive average marginal gain across coalition contexts,
including contexts in which some competing representations are absent. The
result is consistent with contextual substitutability: a representation may
add predictive information when some competitors are absent while receiving
no allocation in the complete ten-representation mixture. This does not imply
that the representation should be retained for deployment.

The result also limits what can be inferred from a zero simplex-gate
coefficient. Full-gate allocation, leave-one-out loss, and Shapley utility
measure different properties of the same predictive resources. A zero
coefficient is therefore insufficient evidence that a representation is
globally irrelevant. Positive Shapley utility likewise does not establish
unique information or causal acoustic content. Mean pooling, linear ridge
experts, and simplex fusion constrain the primary result; whether the same
allocation--utility mismatch holds for nonlinear fusion remains untested.

To test sensitivity to the simplex constraint, we examined the exploratory
positive-NNLS games. In this rule the fitted positive coefficients do not sum
to one (the observed corpus totals ranged from 1.09 to 1.27). Among non-simplex
NNLS gates with zero full-model coefficients, test Shapley utility was positive
in 8/9 cases across BVCC,
SingMOS, and TMHINT-QI, but in 0/4 cases for BRSpeechMOS. NNLS test MSE on
BRSpeechMOS was .904, versus .704 for simplex fusion, consistent with severe
overfitting on its six training systems. Because this check uses test-set
point estimates without validation-frozen discovery or uncertainty intervals,
it is suggestive rather than confirmatory. It supports a limited sensitivity
claim beyond the simplex constraint, not a general claim about arbitrary or
nonlinear fusion. Ridge stacking did not yield exact-zero bottom-ranked
coefficients in most corpora, and the MLP check did not compute coalition
Shapley values.

\textbf{Sensitivity to the player definition and prediction geometry.}
Finally, we test whether the finding depends on the player definition or on
prediction geometry. The original game contained 16 discordant cases. Across alternative player
definitions, the count ranged from 7 to 13 cases, while the qualitative
ranking remained stable (Table~\ref{tab:player_sensitivity}). Grouped-Shapley
rank correlations ranged from $0.976$ to $1.000$, while Owen-value rank
correlations were approximately $1.000$, with sign agreement from $0.889$ to
$1.000$.

\begin{table}[t]
\centering
\small
\caption{Sensitivity of the zero-gate/positive-Shapley rate to the player
definition. Entries report discordant cases divided by eligible
corpus--player pairs. Grouped $F_0$ merges continuous and quantized RMVPE;
grouped prosody merges both RMVPE views with eGeMAPS. Removal rows delete the
indicated players. Rates are descriptive diagnostics, not population
frequencies.}
\label{tab:player_sensitivity}
\setlength{\tabcolsep}{3.0pt}
\begin{tabular}{lrrr}
\toprule
Player definition & Players & Cases/eligible & Rate \\
\midrule
Original                & 10 & $16/40$ & $40.0\%$ \\
Grouped $F_0$ views     & 9  & $13/36$ & $36.1\%$ \\
Grouped prosody         & 8  & $9/32$  & $28.1\%$ \\
Without quantized $F_0$ & 9  & $13/36$ & $36.1\%$ \\
Without eGeMAPS         & 9  & $10/36$ & $27.8\%$ \\
Without both            & 8  & $7/32$  & $21.9\%$ \\
\bottomrule
\end{tabular}
\end{table}

QAP analysis associated prediction correlation with Shapley interaction
strength. Across corpora, the correlation--interaction associations were
$\rho=.368/.943/.443/.957$, with raw
$p=.0069/<10^{-5}/.124/<10^{-5}$. After Holm correction across the four
exploratory tests, the adjusted $p$-values were $.014$, $<.001$, $.124$, and
$<.001$, respectively. Accuracy-matched Gaussian controls yielded
approximately $.30/.96/.80/.89$, indicating that correlated prediction
geometry can reproduce part of the observed interaction pattern without
establishing shared acoustic semantics.

Gate support was less stable than the contextual rankings. Across 100
system-level refits, the four corpora produced 51, 4, 17, and 14 distinct gate
supports, respectively. Full-refit controls produced no discordant cases in 40
diagnostic games. Given the limited null controls, the result is qualitative:
the allocation--utility mismatch persisted across the audited player
definitions, but its count depended on modeling choices.

For a separate numerical check of the complete ten-player gates, we fit each
corpus from a uniform start and 100 symmetric Dirichlet starts
($\mathrm{Dirichlet}(1,\ldots,1)$; seed 20260915). All 101/101 SLSQP runs per
corpus reported success and recovered the same support; the largest loss
difference from the uniform-start solution was $1.05\times10^{-12}$. This
multi-start audit covers the four full gates, not every coalition fit.

\balance
\textbf{Scope and limitations.}
The outer benchmark measures transfer to previously unseen systems within each
corpus, not to unseen corpora. Leave-one-corpus-out results reaggregate the same
held-system errors; they are not models trained on three corpora and tested on
the fourth. The nested benchmark uses one deterministic outer/inner GroupKFold
assignment; its system bootstrap does not measure split-to-split variation, and
we did not pre-specify an equivalence margin or run an equivalence test. The
primary attribution claim is limited to ten frozen linear ridge experts under
a global non-negative simplex gate. Positive NNLS gives an exploratory
test-set sensitivity in three corpora, but not in six-system BRSpeechMOS;
we did not compute a coalition-Shapley audit for nonlinear MLP fusion. Both
inner selection and outer evaluation use equal-system MSE. Finally, checkpoint
hub revisions, weight hashes, and the torchaudio/openSMILE package versions
were not archived, precluding bit-identical embedding re-extraction.

Within global simplex fusion, the results distinguish allocation, contextual
contribution, and deployment value. Exact Shapley analysis identified
contextual utility not visible in full-model weights, but this attribution
did not yield a statistically detectable selection advantage over singleton
ranking on held-out systems. The alternative-fusion NNLS result remains an
exploratory, test-set sensitivity rather than a second confirmatory finding.
For deployment-oriented selection, system-disjoint outer evaluation remains
the relevant criterion, while coalition analysis serves as a diagnostic of
substitutability and representation interaction.


\section{Conclusion}
\label{sec:conclusion}

In this paper, we separated full-mixture allocation, contextual utility, and
fixed-budget deployment value for frozen ridge experts under global simplex
fusion. In the primary audit, 16 of 40 corpus--representation pairs had
full-gate weight below $10^{-6}$ but positive point-estimate Shapley utility;
an exploratory positive-NNLS reanalysis showed the pattern in 8/9 zero-weight
cases across BVCC, SingMOS, and TMHINT-QI, but in 0/4 BRSpeechMOS cases. Exact
Shapley ranking provided no multiplicity-adjusted held-system selection
advantage. Its macro point estimate changed from $-2.15\%$ to $+0.63\%$ when
BRSpeechMOS was excluded, so the ranking is not corpus-robust. These results do
not establish equivalence between selectors or generalize to nonlinear fusion.

These findings distinguish allocation from contextual contribution under
global simplex fusion, but Shapley utility did not replace direct evaluation
of fixed-budget subsets. Selection should be tested under the intended system
shift; repeated group splits and pre-specified equivalence tests are important
next steps. Broader work should also test nonlinear fusion, exact checkpoint
revisions, and selection budgets defined by latency, memory, or robustness.
\section{Acknowledgments}

This work has been partially funded by the project Research and
Development of Algorithms for Construction of Digital Human Technological
Components supported by Advanced Knowledge Center in Immersive Technologies
(AKCIT), with financial resources from the PPI IoT of the MCTI grant number
057/2023, signed with EMBRAPII.

\textit{Generative AI use disclosure:} OpenAI ChatGPT was used only for
language editing, including improvements to grammar, clarity, concision, and
phrasing. The authors take full responsibility for the manuscript.



\bibliographystyle{IEEEbib}
\bibliography{refs}



\end{document}
